\section{Related work and positioning}
Prediction-powered inference corrects an imputed estimate with labeled evidence about prediction error, allowing a pretrained black-box predictor to support valid inference \citep{angelopoulos2023ppi}. Its framework includes both normal-approximation procedures and constructions based on finite-sample concentration; calling an interval ``PPI'' alone therefore does not specify its guarantee. Our unit-control baseline is the usual Gaussian mean interval, while the empirical Bernstein comparisons separately address finite-sample validity.

PPI++ introduces power tuning and asymptotically normal estimators for a broad class of targets \citep{angelopoulos2024ppipp}. For means, power tuning learns a scalar control coefficient. Its discussion also explicitly connects prediction-powered estimators to augmented inverse-propensity estimation and the conditional-mean nuisance function, including the $1/(1+n/N)$ shrinkage factor. We use this established connection rather than claim a new semiparametric principle. Our narrower question concerns what a fixed low-dimensional nonlinear transformation of an already frozen scalar score can add beyond scalar tuning, after paying for fitting and checking coverage.

Cross-prediction-powered inference uses labeled folds to train prediction models while evaluating their corrections on held-out data \citep{zrnic2024cross}. Here the original prediction model remains fixed: cross-fitting applies only to the small inferential regression on its scalar output. We consequently make no claim of introducing cross-fitting or of evaluating the full cross-prediction method with newly trained predictors. The empirical Bernstein construction uses the known-range inequality of \citet{maurer2009empirical}; its guarantee comes from independent training and concentration, rather than from a Gaussian approximation. The study's contribution is the controlled comparison, explicit conditions, and reproducible evidence at this fixed-input scope.

Stratified PPI exploits conditional heterogeneity through a prespecified partition, fixed within-stratum sample counts, and masses treated as known \citep{fisch2024stratified}. Retrospective iid draws instead have random stratum counts and finite-$N$ mass uncertainty; importing its allocation benefit or treating estimated masses as fixed would change this comparison. We do not evaluate an iid poststratified/bin-indicator control, and claim no superiority over it or all nonlinear controls. Active statistical inference improves precision by choosing which labels to acquire \citep{zrnic2024active}; label identities are already fixed here, so that acquisition advantage is outside our setting.
